\scanpage{13}
\margin{12}

- 10 -\hfill n° 234

5. \uline{Matrice d'une forme sesquilinéaire.}

Supposons que $E$ admette une base finie $(e_i)_{1\leq i\leq n}$, et posons
$\rho_{ij}=\Phi(e_i,e_j)$. On a
\[
(15)\qquad
\Phi\left(\sum_{i=1}^n\xi_i e_i,\sum_{i=1}^n\eta_i e_i\right)
=\sum_{i,j}\xi_i\rho_{ij}\eta_j^J.
\]
Inversement, soit $(\rho_{ij})_{1\leq i\leq n,\ 1\leq j\leq n}$ une
\deleted{famille d'}\inserted{matrice à} éléments \deleted{de}\inserted{dans} $A$
et définissons $\Phi$ par la formule (15). Il est immédiat que $\Phi$ est une forme sesquilinéaire sur $E$.
La matrice carrée $\mathcal R=(\rho_{ij})$ à éléments dans $A$, est appelée
\uline{la matrice de $\Phi$ par rapport à la base} $(e_i)$.

Dans tout ce qui suit, pour toute matrice (rectangulaire) $\mathcal U=(\lambda_{ij})$ à éléments dans $A$, nous poserons $\mathcal U^J=(\lambda_{ij}^J)$.
Si on désigne par $x$ et $y$ les matrices à une colonne ayant comme éléments respectifs les $\xi_i$ et les $\eta_i$ $(1\leq i\leq n)$, la formule (15) peut encore s'écrire
\[
(16)\qquad \Phi(x,y)={}^t x.\mathcal R.y^J.
\]

Soient $F$ un second $A$-module à gauche admettant une base finie $(f_j)_{1\leq j\leq m}$,
$u$ une application linéaire de $F$ dans $E$, et posons
$u(f_j)=\sum_{i=1}^n\lambda_{ij}e_i$. Soit $\Phi'$ la transformée de $\Phi$ par $u$, et posons
$\rho'_{ij}=\Phi'(f_i,f_j)$ $(1\leq i\leq m,1\leq j\leq m)$. On a
\[
\rho'_{ij}=\Phi\left(\sum_{k=1}^n\lambda_{ki}e_k,\sum_{k=1}^n\lambda_{kj}e_k\right)
=\sum_{h,k}\lambda_{hi}\rho_{hk}\lambda_{kj}^J
\]
de sorte que, si $\mathcal U$ est la matrice $(\lambda_{ij})$ à $\unsure{m|n}$ lignes et $\unsure{n|m}$ colonnes,
la matrice $\mathcal R'=(\rho'_{ij})$ de $\Phi'$ par rapport à $(f_j)$ est donnée par
\[
(17)\qquad \mathcal R'={}^t\mathcal U.\mathcal R.\mathcal U^J.
\]

Dans cette formule, les matrices ${}^t\mathcal U$ et $\mathcal U^J$ doivent être considérées comme ayant leurs éléments dans $A$ (chap.II,\S6,n°6) ; il en est de même les matrices ${}^t x$ et $y^J$ dans (16).
